% Part: first-order-logic
% Chapter: model-theory
% Section: partial-iso

\documentclass[../../../include/open-logic-section]{subfiles}

\begin{document}

\olfileid{mod}{bas}{dlo}
\section{సాంద్ర రేఖీయ క్రమాలు}

\begin{defn}
  ఒకే ఒక్క ద్విస్థాన \tetoken{విధేయ సంకేతం}{!!{predicate}}~$<$ గల
  \tetoken{భాష}{!!{language}}కు చెందిన, కింది వాక్యాలను సంతృప్తిపరిచే
  \tetoken{నిర్మాణం}{!!a{structure}}~$\Struct{M}$ను
  \emph{అంత్యబిందువులు లేని సాంద్ర రేఖీయ క్రమం} అంటాం:
  \begin{enumerate}
  \item $\lforall[x][\lnot x < x]$;
  \item $\lforall[x][\lforall[y][\lforall[z][(x < y \lif (y < z \lif x
    <z ))]]]$;
  \item $\lforall[x][\lforall[y][(x< y \lor \eq[x][y] \lor y < x)]]$;
  \item $\lforall[x][\lexists[y][x < y]]$;
  \item $\lforall[x][\lexists[y][y < x]]$;
  \item $\lforall[x][\lforall[y][(x < y \lif \lexists[z][(x < z \land
        z < y))]]]$.
  \end{enumerate}
\end{defn}

\begin{thm}\ollabel{thm:cantorQ}
  అంత్యబిందువులు లేని ఏ రెండు \tetoken{లెక్కించదగిన}{!!{enumerable}} సాంద్ర
  రేఖీయ క్రమాలైనా సమరూపాలు.
\end{thm}

\begin{proof}
  $\Struct{M_1}$, $\Struct{M_2}$లు అంత్యబిందువులు లేని
  \tetoken{లెక్కించదగిన}{!!{enumerable}} సాంద్ర రేఖీయ క్రమాలు అనుకుందాం;
  ${<_1}=\Assign{<}{M_1}$, ${<_2}=\Assign{<}{M_2}$గా ఉంచండి. వాటి మధ్య
  అన్ని పాక్షిక సమరూపతల సమితి $\PIso{I}$ అనుకుందాం. కనీసం
  $\emptyset\in\PIso{I}$ కాబట్టి $\PIso{I}$ ఖాళీ కాదు.
  $\PIso{I}$ ముందుకు-వెనుకకు ధర్మాన్ని నెరవేరుస్తుందని చూపుతాం. అప్పుడు
  $\Struct{M_1}\iso[p]\Struct{M_2}$; \olref[pis]{thm:p-isom1} వల్ల
  సిద్ధాంతం వస్తుంది.

  $\PIso{I}$ ముందుకు ధర్మాన్ని నెరవేరుస్తుందని చూపడానికి
  $p\in\PIso{I}$, $a\in\Domain{M_1}$ తీసుకోండి. $p=\emptyset$ అయితే,
  ఏ $b\in\Domain{M_2}$నైనా ఎంచుకొని
  $q=\{\langle a,b\rangle\}$గా ఉంచవచ్చు. ఇక $p\ne\emptyset$
  అనుకుందాం; $i=1$, \dots,~$n$కు $p(a_i)=b_i$గా వ్రాసి, సాధారణతకు
  భంగం లేకుండా $a_1<_1a_2<_1\cdots<_1a_n$ అనుకుందాం.
  $a$ ఇప్పటికే $p$ నిర్వచన క్షేత్రంలో ఉంటే $q=p$గా తీసుకోండి.
  లేదంటే $b\in\Domain{M_2}$ను కింది విధంగా కనుగొనండి:
  \begin{enumerate}
  \item $a<_1a_1$ అయితే, $b<_2b_1$ అయ్యే
    $b\in\Domain{M_2}$ను తీసుకోండి;
  \item $a_n<_1a$ అయితే, $b_n<_2b$ అయ్యే
    $b\in\Domain{M_2}$ను తీసుకోండి;
  \item ఏదో $i$కు $a_i<_1a<_1a_{i+1}$ అయితే,
    $b_i<_2b<_2b_{i+1}$ అయ్యే $b\in\Domain{M_2}$ను తీసుకోండి.
  \end{enumerate}
  \sourcecorrection{OLTEMODBAS-009}{ఆంగ్ల నిరూపణలోని మూడు స్థాన-సందర్భాలు $p$ ఖాళీగా ఉన్న సందర్భాన్నీ, $a$ ఇప్పటికే $p$ నిర్వచన క్షేత్రంలో ఉన్న సందర్భాన్నీ కప్పవు. మొదటిదానికి ఒక యథేచ్ఛ ఏక-జత పాక్షిక సమరూపతను, రెండోదానికి $q=p$ను స్పష్టం చేశాం; కొత్త $a$ కోసం మూలంలోని మూడు క్రమ-సందర్భాలను మార్చలేదు.}
  $\Struct{M_2}$ అంత్యబిందువులు లేని సాంద్ర రేఖీయ క్రమం కాబట్టి,
  కోరిన ధర్మం గల $b$ను ఎల్లప్పుడూ కనుగొనవచ్చు.
  $q=p\cup\{\langle a,b\rangle\}$గా నిర్వచిస్తే,
  $q\in\PIso{I}$ అనేది కోరిన $p$ విస్తరణ. దీనితో ముందుకు ధర్మం
  నిర్ధారితమైంది. వెనుకకు ధర్మం నిరూపణ కూడా ఇదే విధంగా ఉంటుంది.
  కాబట్టి $\Struct{M_1}\iso[p]\Struct{M_2}$; అప్పుడు
  \olref[pis]{thm:p-isom1} వల్ల
  $\Struct{M_1}\iso\Struct{M_2}$.
\end{proof}

\begin{prob}
  $\PIso{I}$ వెనుకకు ధర్మాన్ని నెరవేరుస్తుందని ధ్రువీకరించి
  \olref[mod][bas][dlo]{thm:cantorQ} నిరూపణను పూర్తి చేయండి.
\end{prob}

\begin{rem}
  $\Struct{S}$ అనేది అంత్యబిందువులు లేని ఏ
  \tetoken{లెక్కించదగిన}{!!{enumerable}} సాంద్ర రేఖీయ క్రమమైనా కావచ్చు.
  అప్పుడు (\olref{thm:cantorQ} వల్ల)
  $\Struct{S}\iso\Struct{Q}$; ఇక్కడ
  $\Struct{Q}=(\Rat,<)$ అనేది కరణీయ సంఖ్యల సమితి $\Rat$ను వ్యక్తి
  క్షేత్రంగా గల \tetoken{లెక్కించదగిన}{!!{enumerable}} సాంద్ర రేఖీయ క్రమం.
  ఇప్పుడు \olref[thm]{remark:R}లోని
  \tetoken{నిర్మాణం}{!!{structure}}~$\Struct{R}=(\Real,<)$ను మళ్లీ
  పరిశీలించండి. $\Struct{R}\elemequiv\Struct{S}$ అయ్యే
  \tetoken{లెక్కించదగిన}{!!a{enumerable}}
  \tetoken{నిర్మాణం}{!!{structure}}~$\Struct{S}$ ఉందని చూశాం. కానీ
  $\Struct{S}$ అంత్యబిందువులు లేని
  \tetoken{లెక్కించదగిన}{!!a{enumerable}} సాంద్ర రేఖీయ క్రమం కాబట్టి,
  అది \tetoken{నిర్మాణం}{!!{structure}}~$\Struct{Q}$కు సమరూపం; అందువల్ల
  మొదటిస్థాయి-వాక్య తుల్యం కూడా. మొదటిస్థాయి-వాక్య తుల్యత యొక్క
  సంక్రమణశీలత వల్ల $\Struct{R}\elemequiv\Struct{Q}$.
  (అదే ముందుకు-వెనుకకు వాదనతో
  $\Struct{R}\iso[p]\Struct{Q}$ను స్థాపించి దీన్ని నేరుగా కూడా
  చూపవచ్చు.)
\end{rem}
\end{document}
